% ~500 words. CLIF schema, MIMIC-CLIF entry point, consortium cohort,
% privacy-preserving access mechanisms.
\section{Datasets}
\label{sec:datasets}

\subsection*{The Common Longitudinal ICU Format}

CLIF is a critical-care minimum dataset that harmonizes ICU
electronic health record content into a small set of relational tables
with standardized vocabularies for vital sign and laboratory
categories, medication routes, and respiratory support modes
\citep{rojas2025clif}. As of the founding paper, CLIF spanned
\est{$\sim$9} health systems, \est{$\sim$39} hospitals, and
\est{$\sim$111{,}000} ICU admissions; the consortium has since grown
to 12+ academic medical centers covering \est{$\sim$2{,}000} ICU beds.
The schema is intentionally narrow --- vitals, labs, medications,
respiratory support, admission--discharge--transfer, and patient
demographics --- so that contributing sites can map their EHR exports
once and reuse the mapping across studies. In spirit, CLIF is to
critical care what mCODE \citep{mcode2024hl7} is to oncology: a
clinically scoped common-data-element layer that makes downstream
analytic code portable. CLIF is complementary to broader EHR standards
such as the Medical Event Data Standard \citep{arnrich2025meds} and
OMOP; ETL adapters between CLIF and these formats are maintained so
that adoption does not require schema lock-in.

\subsection*{MIMIC-CLIF as the public development entry point}

To support open development, the consortium distributes MIMIC-CLIF:
the MIMIC-IV critical-care dataset \citep{johnson2023mimiciv,johnson2016mimiciii}
re-expressed in the CLIF v2.1 schema, with \est{$\sim$70{,}000} ICU
stays. Method developers train and self-evaluate against MIMIC-CLIF;
because the schema is identical to consortium-side data, the same code
runs unmodified when later executed federated. This mirrors the
``develop-on-public, validate-on-private'' pattern that has become
canonical in clinical AI but is rarely operationalized end-to-end.

\subsection*{Consortium cohort}

\paragraph{Coverage.} Across the 12+ sites currently in CLIF, FLAIR
operates on an aggregate cohort of \est{$\sim$500{,}000} ICU stays
spanning urban academic medical centers, large community hospitals, and
health systems with mixed payer populations. Demographic distribution
(sex, race, ethnicity, age), admission SOFA, and Charlson and
Elixhauser comorbidity indices are summarized per task in the cohort
report (\texttt{table1.json}); a pooled Table~1 with HIPAA-compliant
cell suppression accompanies the manuscript. % TODO-NUMBERS

\paragraph{Privacy-preserving access.} Three properties together
constitute FLAIR's access model. \emph{No patient-level data leaves
any site}: outputs are JSON aggregates produced by the FLAIR library
inside the container. \emph{Library-enforced suppression}: counts
below 10 in cohort tables become \texttt{"<10"} in
\texttt{table1.json}; demographic subgroups with $n<10$ are dropped
entirely from \texttt{fairness.json}; calibration bins with fewer than
10 observations report a null event rate in \texttt{curves.json}.
\emph{Verifiable execution}: a SHA-256 digest of the task script and
the FLAIR library version are stamped into every output JSON, so the
central aggregator can confirm that the same code ran byte-for-byte at
each site. Network isolation is enforced one layer below, by the
CLIF-WorkBench container in which the library runs --- no internet,
no root, immutable read-only image --- consistent with the spirit of
analysis platforms such as OpenSAFELY \citep{opensafely2021naturemedicine}
in which the code travels and the data does not. Together these
properties move evidence generation from the slow, bilateral DUA path
to a routine, repeatable operation.
